\bta{2012}

\begin{enumerate}
	%\renewcommand{\labelenumi}{\arabic{enumi}.}
	% A(\Alph) a(\alph) I(\Roman) i(\roman) 1(\arabic)
	%设定全局标号series=example	%引用全局变量resume=example
	%[topsep=-0.3em,parsep=-0.3em,itemsep=-0.3em,partopsep=-0.3em]
	%可使用leftmargin调整列表环境左边的空白长度 [leftmargin=0em]
	\item
	（共 30 分）
	质量为 $\mu$ 的粒子在一维无限深势阱中运动, 势能
	$$
	V=\left\{\begin{aligned}
		&0, & 0 \leq x \leq a \\
		&\infty, & x<0 \text { 或 } x>a
	\end{aligned}\right.
	$$
	\begin{enumerate}
		%\renewcommand{\labelenumi}{\arabic{enumi}.}
		% A(\Alph) a(\alph) I(\Roman) i(\roman) 1(\arabic)
		%设定全局标号series=example	%引用全局变量resume=example
		%[topsep=-0.3em,parsep=-0.3em,itemsep=-0.3em,partopsep=-0.3em]
		%可使用leftmargin调整列表环境左边的空白长度 [leftmargin=0em]
		\item
		求粒子的能级和归一化波函数。
		\item
		画出处于第二、第三激发态的粒子的概率密度的示意图。
		\item
		求坐标算符在能量表象下的矩阵元。
		
		
		
	\end{enumerate}
\banswer{
\begin{enumerate}
\item 阱内方程为 $\psi''+k^2\psi=0$，其中 $k^2=2\mu E/\hbar^2$。
通解 $A\sin kx+B\cos kx$ 由 $\psi(0)=0$ 得 $B=0$，由 $\psi(a)=0$ 得
$k=n\pi/a$。再由 $|A|^2\int_0^a\sin^2(n\pi x/a)dx=1$ 得
\[
 \boxed{\psi_n(x)=\sqrt{\frac2a}\sin\frac{n\pi x}{a},\quad
 E_n=\frac{n^2\pi^2\hbar^2}{2\mu a^2}},\quad n=1,2,\ldots,
\]
阱外波函数为零。
\item 第二、第三激发态分别为 $n=3,4$，概率密度为
\[
 \rho_3(x)=\frac2a\sin^2\frac{3\pi x}{a},\qquad
 \rho_4(x)=\frac2a\sin^2\frac{4\pi x}{a}.
\]
对 $\rho_n$，零点为 $x=ja/n$（$j=0,\ldots,n$），极大值位于
$x=(2j+1)a/(2n)$（$j=0,\ldots,n-1$），高度均为 $2/a$。
因此两图分别有三个和四个等高波瓣。
\item 对角元由关于 $a/2$ 的反射对称性得 $x_{nn}=a/2$。
当 $n\ne m$ 时，用 $2\sin A\sin B=\cos(A-B)-\cos(A+B)$，
\begin{align*}
 x_{nm}&=\frac2a\int_0^a x\sin\frac{n\pi x}{a}\sin\frac{m\pi x}{a}dx\\
 &=\frac1a\int_0^a x\left[\cos\frac{(n-m)\pi x}{a}
                         -\cos\frac{(n+m)\pi x}{a}\right]dx.
\end{align*}
对非零整数 $j$，
\[
 \int_0^a x\cos\frac{j\pi x}{a}dx
 =\left[\frac{ax}{j\pi}\sin\frac{j\pi x}{a}
       +\frac{a^2}{j^2\pi^2}\cos\frac{j\pi x}{a}\right]_0^a
 =\frac{a^2[(-1)^j-1]}{j^2\pi^2}.
\]
故
\[
 \boxed{x_{nm}=\begin{cases}
 a/2,&n=m,\\[2pt]
 \dfrac{a}{\pi^2}\left[\dfrac{(-1)^{n-m}-1}{(n-m)^2}
                       -\dfrac{(-1)^{n+m}-1}{(n+m)^2}\right],&n\ne m.
 \end{cases}}
\]
特别地，$n\ne m$ 同奇偶时为零，异奇偶时为
$-8anm/[\pi^2(n^2-m^2)^2]$。
\end{enumerate}
}
	
\item 
（共 30 分）
质量为 $\mu$ 的一维谐振子, 带电量为 $q$， 初始 $t=-\infty$ 时处于基态 $ \ket{0}$。 设
加上微扰 $H^{\prime}=-q E x \exp \left(-\frac{t^{2}}{\tau^{2}}\right)$, 其中 $E$ 为外电场强度， $\tau$ 为参数。求 $t=+\infty$ 时 谐振子仍停留在基态的概率。

\banswer{
令 $x_0=\sqrt{\hbar/(2\mu\omega)}$，则 $x=x_0(a+a^\dagger)$，从基态出发
\[
 \langle n|H'(t)|0\rangle=-qEx_0e^{-t^2/\tau^2}\delta_{n1}.
\]
因此在一阶振幅近似下，只有 $n=1$ 被激发：
\begin{align*}
 c_1^{(1)}(\infty)
 &=-\frac{i}{\hbar}\int_{-\infty}^\infty
           (-qEx_0)e^{-t^2/\tau^2}e^{i\omega t}dt\\
 &=\frac{iqEx_0}{\hbar}\sqrt\pi\tau e^{-\omega^2\tau^2/4}.
\end{align*}
这里使用 Gaussian Fourier 积分
$\int e^{-t^2/\tau^2+i\omega t}dt=\sqrt\pi\tau e^{-\omega^2\tau^2/4}$。
令
\[
 \eta=|c_1^{(1)}|^2
     =\frac{\pi q^2E^2\tau^2}{2\mu\hbar\omega}e^{-\omega^2\tau^2/2}.
\]
由总概率守恒，在弱场 $\eta\ll1$ 时，
\[
 \boxed{P_0=1-\eta+O(E^4)}.
\]
不能把“只有第一激发态”理解成任意场强下的精确结论。线性驱动谐振子的精确末态是相干态，
其参数为 $\alpha=iqEx_0\sqrt\pi\tau e^{-\omega^2\tau^2/4}/\hbar$，故精确地有
$P_0=e^{-|\alpha|^2}=e^{-\eta}$；展开后与微扰结果一致。
}

\item 
（共 30 分）
设一转动惯量为 $I$ 、电偶极矩为 $D$ 的转子自由地在 $X-Y$ 平面内
转动， 相对$ X $轴正向的转角为 $\varphi$ 。
\begin{enumerate}
	%\renewcommand{\labelenumi}{\arabic{enumi}.}
	% A(\Alph) a(\alph) I(\Roman) i(\roman) 1(\arabic)
	%设定全局标号series=example	%引用全局变量resume=example
	%[topsep=-0.3em,parsep=-0.3em,itemsep=-0.3em,partopsep=-0.3em]
	%可使用leftmargin调整列表环境左边的空白长度 [leftmargin=0em]
	\item
试求其能量本征值和本征态。
\item 
设沿 $\mathrm{X}$ 方向加上电场 $ \mathscr{E}$, 即微扰哈密顿量为 $H^{\prime}=-D \mathscr{E} \cos \varphi$,	试用微扰论求其
基态能量的一级和二级微扰修正。
\end{enumerate}
\banswer{
\begin{enumerate}
\item 转子的自由 Hamilton 量为 $H_0=-\hbar^2\partial_\varphi^2/(2I)$。
通解中的 $e^{im\varphi}$ 必须满足 $\psi(\varphi+2\pi)=\psi(\varphi)$，故
$m\in\mathbb Z$。由 $\int_0^{2\pi}|\psi|^2d\varphi=1$，得
\[
 \boxed{\psi_m(\varphi)=\frac{e^{im\varphi}}{\sqrt{2\pi}},\quad
 E_m^{(0)}=\frac{m^2\hbar^2}{2I}}.
\]
$m=0$ 为非简并基态，$m\ne0$ 与 $-m$ 简并。
\item 基态为常函数，所以
\[
 E_0^{(1)}=-\frac{D\mathscr E}{2\pi}\int_0^{2\pi}\cos\varphi\,d\varphi=0.
\]
又由 $\cos\varphi=(e^{i\varphi}+e^{-i\varphi})/2$ 及指数函数正交性，
\[
 H'_{m0}=-\frac{D\mathscr E}{2}
           (\delta_{m,1}+\delta_{m,-1}).
\]
两个非零矩阵元均为 $-D\mathscr E/2$，能量分母均为 $-\hbar^2/(2I)$，故
\[
 \boxed{E_0^{(2)}=2\frac{(D\mathscr E/2)^2}{-\hbar^2/(2I)}
                  =-\frac{D^2\mathscr E^2I}{\hbar^2}}.
\]
因子 $2$ 来自 $m=+1,-1$ 两个中间态；适用条件为 $|D\mathscr E|I/\hbar^2\ll1$。
\end{enumerate}
}


\newpage
\item 
（共 30 分）
\begin{enumerate}
	%\renewcommand{\labelenumi}{\arabic{enumi}.}
	% A(\Alph) a(\alph) I(\Roman) i(\roman) 1(\arabic)
	%设定全局标号series=example	%引用全局变量resume=example
	%[topsep=-0.3em,parsep=-0.3em,itemsep=-0.3em,partopsep=-0.3em]
	%可使用leftmargin调整列表环境左边的空白长度 [leftmargin=0em]
	\item
写出角动量算符的三个分量 $J_{x} $、$ J_{y} $、$ J_{z}$ 相互间满足的所有对易关系。
	\item 
	试利用这些对易关系, 证明矩阵元 $\left\langle m\left|J_{x}\right| n\right\rangle$ 仅当 $m=n \pm 1$ 时不为零。其中 $|m\rangle$ 、 $|n\rangle$ 分别为 $J_{z}$ 的本征值为 $m \hbar $、$ n \hbar$ 的本征态。
	
	\item 
	设角动量量子数 $j=1$ 。 已知在 $J_{z}$ 的某一个本征态 $|m\rangle$ 中, $J_{x}$ 取值为 0 的概 率为 $1 / 2$。求 $J_{x}$ 取值为 $-\hbar$ 的概率。
	
\end{enumerate}

\banswer{
\begin{enumerate}
	%\renewcommand{\labelenumi}{\arabic{enumi}.}
	% A(\Alph) a(\alph) I(\Roman) i(\roman) 1(\arabic)
	%设定全局标号series=example	%引用全局变量resume=example
	%[topsep=-0.3em,parsep=-0.3em,itemsep=-0.3em,partopsep=-0.3em]
	%可使用leftmargin调整列表环境左边的空白长度 [leftmargin=0em]
		\item
		三个非平凡对易关系及其反序形式为
		\[
			\boxed{[J_x,J_y]=i\hbar J_z,\qquad
			[J_y,J_z]=i\hbar J_x,\qquad
			[J_z,J_x]=i\hbar J_y},
		\]
		\[
			[J_y,J_x]=-i\hbar J_z,\quad
			[J_z,J_y]=-i\hbar J_x,\quad
			[J_x,J_z]=-i\hbar J_y,
		\]
		且 $[J_i,J_i]=0$。
		\item
		由 $[J_z,J_x]=i\hbar J_y$ 和
		$[J_z,J_y]=-i\hbar J_x$，
		\[
			(m-n)\langle m|J_x|n\rangle
			=i\langle m|J_y|n\rangle,
		\]
		\[
			(m-n)\langle m|J_y|n\rangle
			=-i\langle m|J_x|n\rangle.
		\]
		消去 $J_y$ 的矩阵元可得
		\[
			\left[(m-n)^2-1\right]\langle m|J_x|n\rangle=0.
		\]
		故 $\langle m|J_x|n\rangle$ 仅在
		$\boxed{m=n\pm1}$ 时可能不为零。
	\item 
	由上一问的选择定则，$\langle m|J_x|m\rangle=0$。
	$J_x$ 的三个测值为 $\hbar,0,-\hbar$，记其概率为 $P_+,P_0,P_-$。
	于是 $\hbar(P_+-P_-)=0$，而题设 $P_0=1/2$、归一化要求
	$P_++P_-=1/2$。两式联立得到
	\[
	 \boxed{P_-=P_+=\frac14}.
	\]
\end{enumerate}


}
\item 
（共 30 分）
氢原子的哈密顿量为 $H_{0}=\frac{p^{2}}{2 \mu}-\frac{e^{2}}{r} $， 基态波函数及基态能量为
$$
\psi_{0}(r)=\frac{1}{\sqrt{\pi a_{0}{ }^{3}}} e^{-\frac{r}{a_{0}}} \qquad  E_{0}=-\frac{e^{2}}{2 a_{0}}=-\frac{\mu e^{4}}{2 \hbar^{2}}
$$
其中 $a_{0}=\frac{\hbar^{2}}{\mu e^{2}}$ 为第一 Bohr 轨道半径。
设体系受到微扰 $H^{\prime}=e \varepsilon z$ 的作用（沿 $z$ 方向加上均匀电场 $\varepsilon$ ), 哈密顿量变成 $H=H_{0}+H^{\prime}$。
\begin{enumerate}
	%\renewcommand{\labelenumi}{\arabic{enumi}.}
	% A(\Alph) a(\alph) I(\Roman) i(\roman) 1(\arabic)
	%设定全局标号series=example	%引用全局变量resume=example
	%[topsep=-0.3em,parsep=-0.3em,itemsep=-0.3em,partopsep=-0.3em]
	%可使用leftmargin调整列表环境左边的空白长度 [leftmargin=0em]
	\item
	计算对易关系: $\left[H_{0}, H^{\prime}\right]$ 及 $\left[H^{\prime},\left[H_{0}, H^{\prime}\right]\right]$ 。
	\item 
	计算 $\psi_{0}$ 下的平均值： $\left\langle H^{\prime}\right\rangle_{0}$ 及 $\left\langle H^{\prime 2}\right\rangle_{0}$ 。
	\item 
	取基态试探波函数为 $\psi(\lambda)=N\left(1+\lambda H^{\prime}\right) \psi_{0}$, 其中 $N$ 为归一化常数。试以 $\lambda$ 为
	变分参数，用变分法求$  H  $的基态能量上限（准确到
$\varepsilon^{2}$量级）。
	
	
\end{enumerate}

\banswer{
\begin{enumerate}
	%\renewcommand{\labelenumi}{\arabic{enumi}.}
	% A(\Alph) a(\alph) I(\Roman) i(\roman) 1(\arabic)
	%设定全局标号series=example	%引用全局变量resume=example
	%[topsep=-0.3em,parsep=-0.3em,itemsep=-0.3em,partopsep=-0.3em]
	%可使用leftmargin调整列表环境左边的空白长度 [leftmargin=0em]
	\item
	势能与 $z$ 对易，且
	\[
	 [p_z^2,z]=p_z[p_z,z]+[p_z,z]p_z=-2i\hbar p_z.
	\]
	故
	\[
	 [H_0,H']=-\frac{i\hbar e\varepsilon}{\mu}p_z,\qquad
	 [H',[H_0,H']]=-\frac{i\hbar(e\varepsilon)^2}{\mu}[z,p_z]
	 =\boxed{\frac{e^2\varepsilon^2\hbar^2}{\mu}}.
	\]
	\begin{note}
		这里直接采用直角坐标系运算，采用球坐标系反而麻烦了。球坐标系下的结果是
		\[ 
		\left[H_{0}, H^{\prime}\right]=-\frac{e \varepsilon \hbar^{2}}{\mu} (\cos \theta \frac{\partial }{\partial r} -\frac{1}{r} \sin \theta \frac{\partial}{\partial \theta} )
		 \]
	\end{note}
	\item 
		基态为偶宇称，$z$ 为奇算符，故 $\langle H'\rangle_0=0$。
	球对称性给出 $\langle z^2\rangle=\langle r^2\rangle/3$，并且
	\[
	 \langle r^2\rangle=\frac4{a_0^3}\int_0^\infty r^4e^{-2r/a_0}dr
	 =\frac4{a_0^3}\frac{4!}{(2/a_0)^5}=3a_0^2.
	\]
	因此 $\boxed{\langle H'^2\rangle_0=e^2\varepsilon^2a_0^2}$。
			\item
		记 $A=H'=e\varepsilon z$，并令
		$s=\langle0|A^2|0\rangle=e^2\varepsilon^2a_0^2$。
		归一化常数满足
		\[
			N^{-2}=1+\lambda^2s.
		\]
		保留到 $\varepsilon^2$，变分能量为
		\[
			E(\lambda)=E_0+2\lambda s+\lambda^2D,
			\qquad
			D=\langle0|A(H_0-E_0)A|0\rangle.
		\]
		利用第一问的二重对易子，
		\[
			2D=\langle0|[A,[H_0,A]]|0\rangle
			=\frac{e^2\varepsilon^2\hbar^2}{\mu},
		\]
		故 $D=e^2\varepsilon^2\hbar^2/(2\mu)$。极小化得到
		\[
			\lambda_{\rm opt}=-\frac{s}{D}
			=-\frac{2\mu a_0^2}{\hbar^2},
		\]
		以及
		\[
			E_{\rm var}=E_0-\frac{s^2}{D}
			=E_0-\frac{2\mu e^2\varepsilon^2a_0^4}{\hbar^2}.
		\]
		再用 $a_0=\hbar^2/(\mu e^2)$，最终
		\[
			\boxed{E_{\rm var}=E_0-2\varepsilon^2a_0^3
			+O(\varepsilon^4)}.
		\]
		负号表示基态在静电场中的二阶 Stark 位移向下。
\end{enumerate}

	
}




\end{enumerate}
